\section{Additional Experimental Results}\label{app_sec:additional_experimental_results}
\subsection{Information Extraction}

\begin{table}
\centering
\resizebox{0.48\textwidth}{!}{
\setlength{\tabcolsep}{5pt}
\begin{tabular}{@{}lccc@{}}
\toprule
& \multicolumn{3}{c}{SynthIE-text} \\
\hspace{4mm} Method & Precision & Recall & F1  \\
\midrule
\midrule
\emph{\textbf{Few-shot Unconstrained}}\\
\hspace{4mm} LLaMA-7B-sc  (4 shots)   &    9.7     &    8.6    &    9.2   \\
\hspace{4mm} LLaMA-13B-sc (4 shots)     &    6.7     &    12.5     &    8.7   \\
\hspace{4mm} LLaMA-33B-sc  (4 shots)    &    10.6     &   20.6     &    14.0   \\
\midrule
\emph{\textbf{Few-shot Constrained}}\\
\hspace{4mm} LLaMA-7B-sc (4 shots)    &   24.1    &    19.3     &    21.5      \\
\hspace{4mm} LLaMA-13B-sc (4 shots)    &    26.4    &    30.5     &    28.3      \\
\hspace{4mm} LLaMA-33B-sc  (4 shots)    &    31.3    &    36.2     &    33.6      \\
\bottomrule
\end{tabular}
}
\caption{\textbf{Results for cIE with subject-collapsed linearization.} }
\label{tab:subject-collapsed-linearization}
\end{table}
\begin{table}
\centering
\resizebox{0.48\textwidth}{!}{
\setlength{\tabcolsep}{5pt}
\begin{tabular}{@{}lccc@{}}
\toprule
& \multicolumn{3}{c}{SynthIE-text} \\
\hspace{4mm} Method & Precision & Recall & F1  \\
\midrule
\midrule
%\emph{\textbf{Micro}} \\
% \hspace{4mm} SotA Pipeline (R)    &    43.30     &    41.73 {\scriptsize± 0.13}    &    42.50 {\scriptsize± 0.13}    &    17.89 {\scriptsize± 0.24} \\
\emph{\textbf{Supervised}} \\
\hspace{4mm} GenIE T5-base-fe~\cite{josifoski-etal-2022-genie}     &    \textbf{49.1}     &    26.7     &    \textbf{34.6 }  \\
\midrule
\emph{\textbf{Few-shot Unconstrained}}\\
\hspace{4mm} LLaMA-7B-feR  (4 shots)   &    9.9     &    13.5     &    11.4   \\
\hspace{4mm} LLaMA-13B-feR (4 shots)     &    10.8     &    17.5     &    13.4   \\
\hspace{4mm} LLaMA-33B-feR  (4 shots)    &    14.0    &    22.5     &    17.3   \\
\hspace{4mm} Vicuna-13B-feR  (4 shots)   &    12.5     &    16.7     &    14.3   \\
\midrule
\emph{\textbf{Few-shot Constrained}}\\
\hspace{4mm} LLaMA-7B-feR (4 shots)   &    28.1     &    23.6     &    25.7   \\
\hspace{4mm} LLaMA-13B-feR (4 shots)     &    31.8     &    31.2     &    31.5   \\
\hspace{4mm} LLaMA-33B-feR  (4 shots)    &    36.4    &    34.8     &    35.6   \\
\hspace{4mm} Vicuna-13B-feR (4 shots)     &    40.3     &    23.8     &    29.9   \\
\bottomrule
\end{tabular}
}
\caption{\textbf{Results for cIE with all relations added to the prompt.} feR stands for fully-expanded linearization with relations added to the prompt.}
\label{tab:InformationExtractionResults_Relations}
\end{table}

\xhdr{Subject-Collapsed Linearization}
Subject-collapsed linearization is a variant of linearization where the subject is collapsed into the object.
It has the advantage of being more compact (token efficient) than the fully-expanded linearization, but it yields slightly lower performance, intuitively because it's less explicit.
Table~\ref{tab:subject-collapsed-linearization} shows that the constrained decoding approach brings a significant improvement over the unconstrained decoding approach for LMs of all sizes.
%
Compared with the results in Table~\ref{tab:subject-collapsed-linearization}, the performance of the subject-collapsed linearization is indeed lower than the fully-expanded linearization.

\xhdr{Adding all relations to the prompt}
With a maximum context length of 2048 tokens, it's not possible to add all the millions of \textbf{entities} to the prompt.
However, we can add all the \textbf{relations} to the prompt, which is a much smaller set (888 relations).
%
One may wonder if adding all the \textbf{relations} to the prompt would help the model to learn the relation extraction task better.
Because the model can copy the \textbf{relation} from the prompt and this may make the task easier.
However, we find that adding all the \textbf{relations} to the prompt does not help the model to learn the relation extraction task better.
Table~\ref{tab:InformationExtractionResults_Relations} shows that adding all the \textbf{relations} to the prompt only brings a small improvement over the baseline.



\xhdr{Constituency Parsing with Vicuna}
\begin{table}[h]
\centering
\resizebox{0.48\textwidth}{!}{
\setlength{\tabcolsep}{5pt}
\begin{tabular}{@{}lcccc@{}}
\toprule
\hspace{4mm} Method & Precision & Recall & Tag Accuracy & Validity \\
\midrule
\emph{\textbf{Unconstrained}}\\
\hspace{4mm} Vicuna-7B     &    13.5     &    12.7     &    27.8   &    42.2 \\
\hspace{4mm} Vicuna-13B    &    31.6     &    28.1     &    29.6   &    51.4 \\
\midrule
\emph{\textbf{Constrained IIG}}\\
\hspace{4mm} Vicuna-7B     &    17.4     &    16.3     &    30.5   &    54.3 \\
\hspace{4mm} Vicuna-13B    &    30.8     &    27.7     &    33.4   &    56.1 \\
\midrule
\emph{\textbf{Constrained IDG}}\\
\hspace{4mm} Vicuna-7B     &    35.6     &    31.9     &    41.4   &    100.0 \\
\hspace{4mm} Vicuna-13B    &    51.6     &    44.4     &    42.4   &    100.0 \\
\bottomrule
\end{tabular}
}
\caption{Constituency parsing results with Vicuna. The experiments setting is the same as in Table~\ref{tab:ConstituencyParsingResults}.}
\label{tab:ConstituencyParsingResultsWithVicuna}
\end{table}
We also evaluate the constrained decoding approach on the constituency parsing task with the instruction-tuned model Vicuna.
Table~\ref{tab:ConstituencyParsingResultsWithVicuna} shows Vicuna's performance is even worse than the LLaMA model.



